\section{Delivery and control}
\label{sec:delivery}
% Generated by scripts/export_section7.py — do not edit by hand.
\newcommand{\PEPdoaAud}{\$20{,}000\xspace}
\newcommand{\PEPdoaK}{\$20k\xspace}
\newcommand{\PEPdoaHours}{48\xspace}
\newcommand{\PEPcontingency}{\$225{,}000\xspace}
\newcommand{\PEPbase}{\$1{,}500{,}000\xspace}
\newcommand{\PEPbac}{\$1{,}725{,}000\xspace}
\newcommand{\PEPoptionC}{\$21{,}000\xspace}
\newcommand{\PEPpmName}{Abdul Wahab Obeid\xspace}
\newcommand{\PEPmTwoShort}{2~Oct\xspace}
\newcommand{\PEPmFourShort}{15~Oct\xspace}
\newcommand{\PEPmFiveShort}{30~Oct\xspace}
\newcommand{\PEPmSixShort}{10~Nov\xspace}
\newcommand{\PEPmSevenShort}{15~Nov\xspace}
\newcommand{\PEPmFourLong}{15~October\xspace}
\newcommand{\PEPmFiveLong}{30~October\xspace}
\newcommand{\PEPmSevenLong}{15~November\xspace}
\newcommand{\PEPgreenLo}{0.95}
\newcommand{\PEPgreenHi}{1.05}
\newcommand{\PEPamberLo}{0.90}
\newcommand{\PEPredLt}{0.90}
\newcommand{\PEPtZeroSlip}{2\xspace}
\newcommand{\PEPrOiImpact}{\$90{,}000\xspace}
\newcommand{\PEPrOiEmv}{\$27{,}000\xspace}
\newcommand{\PEPrViiImpact}{\$70{,}000\xspace}
\newcommand{\PEPrViiEmv}{\$21{,}000\xspace}
\newcommand{\PEPpvmFour}{\$825{,}000\xspace}
\newcommand{\PEPpvmFournum}{825{,}000}
\newcommand{\PEPbasenum}{1{,}500{,}000}
\newcommand{\PEPbacnum}{1{,}725{,}000}
\newcommand{\PEPonPlanEV}{\$810{,}000\xspace}
\newcommand{\PEPonPlanEVnum}{810{,}000}
\newcommand{\PEPonPlanAC}{\$840{,}000\xspace}
\newcommand{\PEPonPlanACnum}{840{,}000}
\newcommand{\PEPonPlanCPI}{0.96}
\newcommand{\PEPonPlanSPI}{0.98}
\newcommand{\PEPonPlanEac}{\$1{,}558{,}750\xspace}
\newcommand{\PEPonPlanEacnum}{1{,}558{,}750}
\newcommand{\PEPonPlanIeac}{\$1{,}796{,}875\xspace}
\newcommand{\PEPonPlanIeacnum}{1{,}796{,}875}
\newcommand{\PEPlateEV}{\$760{,}000\xspace}
\newcommand{\PEPlateEVnum}{760{,}000}
\newcommand{\PEPlateAC}{\$830{,}000\xspace}
\newcommand{\PEPlateACnum}{830{,}000}
\newcommand{\PEPlateCPI}{0.92}
\newcommand{\PEPlateSPI}{0.92}
\newcommand{\PEPlateEac}{\$1{,}634{,}348\xspace}
\newcommand{\PEPlateEacnum}{1{,}634{,}348}
\newcommand{\PEPmFourGap}{\$15{,}000\xspace}
\newcommand{\PEPpmBox}{%
\textbf{Campaign PM}\\[-0.05em]\textbf{\mbox{\PEPpmName}}\\[0.12em]
  \scriptsize DoA: $\le$ \PEPdoaK{} \emph{and} $\le$ \PEPdoaHours{}~h\\[-0.05em]
  \scriptsize Cannot move M-5 (\PEPmFiveShort) or M-7 (\PEPmSevenShort)}


This section converts the baselines in Sections~\ref{sec:schedule}--\ref{sec:risk} into decision rights, a pad-specific procurement choice, hold points, Bowen HSE, and an earned-value framework that will be read at M-4. It is a control design, not a progress report: the campaign has not been executed.

\subsection{Governance and reporting}
\label{sec:delivery-governance}

Governance is the organisation that will actually stop work, not a generic pyramid \citep{pmi2021,kerzner2017}. The Gilmour Executive Board is the sponsor. Abdul Wahab Obeid is Campaign Project Manager and directs the three consultancy leads in Table~\ref{tab:raci}. Field titles on Figure~\ref{fig:governance} match that table: Launch Director, Range Safety Officer, GSE Lead, Avionics Specialist, Juru Cultural Officer and the design-and-construct package manager.

Delegation is numeric. Abdul may approve a variance of at most \PEPdoaAud \emph{and} at most \PEPdoaHours~hours, and only if M-5 (\PEPmFiveLong, A-133) and M-7 (\PEPmSevenLong, A-143) do not move. Anything else---a contingency draw, Option~C's \PEPoptionC, a hold-point waiver, a $T$-0 change---goes to the Board within 24~hours \citep{kerzner2017}. R-16 (RSO versus Launch Director deadlock) escalates 24~hours before A-141's late finish of \PEPmSixShort, not inside the review room. The \PEPdoaAud ceiling is below the dollar \emph{impact} of the hot-fire rows that would recycle the pad (R-01 \PEPrOiImpact; R-07 \PEPrViiImpact). Residual EMVs on those rows are \PEPrOiEmv and \PEPrViiEmv because treatments already sit in the work baseline; a genuine recycle still cannot be buried in PM discretion.

Cadence matches the physical work. Field leads hold a daily pad standup. After Gate~1 (A-111, \PEPmTwoShort) Ilsa publishes a weekly EVM dashboard from A-112. Abdul takes a fortnightly pack to the sponsor. Daily standup does not rewrite Table~\ref{tab:precedence}. A-112 has 53~days of float so the control system does not compete with stacking for pad occupancy.

\begin{figure}[H]
\centering
\resizebox{\textwidth}{!}{%
\begin{tikzpicture}[
  font=\sffamily,
  >=Stealth,
  board/.style={draw=navy, fill=navy, text=white, rounded corners=2pt,
    inner xsep=8pt, inner ysep=6pt, font=\bfseries\small, align=center, text width=8.4cm},
  pm/.style={draw=navy, fill=teal, text=white, rounded corners=2pt,
    inner xsep=6pt, inner ysep=5pt, font=\small, align=center, text width=7.6cm},
  lead/.style={draw=navy, fill=white, rounded corners=2pt,
    inner xsep=4pt, inner ysep=4pt, font=\scriptsize, align=center,
    text width=3.85cm, minimum height=1.22cm},
  field/.style={draw=navy!70, fill=white, rounded corners=1pt,
    inner xsep=2.5pt, inner ysep=3pt, font=\scriptsize, align=center,
    text width=2.28cm, minimum height=1.08cm},
  note/.style={draw=crimson, dashed, thick, rounded corners=2pt, fill=white,
    inner xsep=6pt, inner ysep=4pt, font=\scriptsize, align=center, text width=15.4cm}
]
\node[board] (board) {Gilmour Executive Board \\[0.12em] \normalfont\scriptsize Sponsor --- contingency, M-5 / M-7, hold-point waivers};
\node[pm] (pm) at (0,-2.05) {\PEPpmBox};
\draw[navy, thick, ->] (board.south) -- (pm.north);

\node[lead] (hadi)  at (-4.55,-4.05) {\textbf{\mbox{Hadi Alkhub}}\\Lead Planner\\schedule, network, compression};
\node[lead] (ziyad) at ( 0.00,-4.05) {\textbf{\mbox{Ziyad Alahmari}}\\Risk and Quality Lead\\register, hold points, HSE};
\node[lead] (ilsa)  at ( 4.55,-4.05) {\textbf{\mbox{Ilsa Hashmi}}\\Cost and Resource Lead\\PMB, EVM, Juru commercial};
\foreach \n in {hadi,ziyad,ilsa}
  \draw[navy, thick, ->] (pm.south) -- ++(0,-0.26) -| (\n.north);

\node[field] (ld)   at (-6.35,-6.05) {\textbf{Launch}\\Director};
\node[field] (rso)  at (-3.81,-6.05) {\textbf{Range Safety}\\Officer};
\node[field] (gse)  at (-1.27,-6.05) {\textbf{GSE}\\Lead};
\node[field] (avio) at ( 1.27,-6.05) {\textbf{Avionics}\\Specialist};
\node[field] (juru) at ( 3.81,-6.05) {\textbf{Juru Cultural}\\Officer};
\node[field] (dc)   at ( 6.35,-6.05) {\textbf{D\&C package}\\manager};
\coordinate (fieldbar) at (0,-5.15);
\draw[navy, thick] (hadi.south) -- ++(0,-0.18) -| (fieldbar);
\draw[navy, thick] (ziyad.south) -- (fieldbar);
\draw[navy, thick] (ilsa.south) -- ++(0,-0.18) -| (fieldbar);
\draw[navy, thick] (-6.35,-5.15) -- (6.35,-5.15);
\foreach \n in {ld,rso,gse,avio,juru,dc}
  \draw[navy!80, ->] (\n.north |- fieldbar) -- (\n.north);

\node[note] (cadence) at (0,-7.35) {\textbf{Reporting cadence.} Daily pad standup \quad$\cdot$\quad weekly EVM dashboard (Ilsa) \quad$\cdot$\quad fortnightly sponsor review \quad$\cdot$\quad 24~h to Board if DoA exceeded or R-16 fires.};
\end{tikzpicture}%
}
\caption[Campaign governance structure]{Campaign governance for Bowen pad operations. Named roles, a numeric DoA on the PM, and the reporting cadence. Field titles match Table~\ref{tab:raci}. The Board, not the PM, moves M-5 or M-7 or draws the \PEPcontingency reserve.}
\label{fig:governance}
\end{figure}

\subsection{Delivery model and procurement}
\label{sec:delivery-model}

The chosen delivery model is a managing-contractor / design-and-construct hybrid \citep{turner2014,pmi2021}. Gilmour retains propulsion intellectual property, vehicle integration and the Australian launch licence. Civil works and cryogenic GSE packages are let as design-and-construct, with an interface-control freeze and liquidated damages. That split is the contractual answer to R-07 and R-08: an ICD lag is an LD event on the contractor, and a latent pad-civil defect stays with the D\&C package after HP-1 NDT (R-08 EMV \$0 to Gilmour). Abdul stays Accountable on that row; the package manager stays Responsible.

Turnkey EPC was considered and rejected. After TestFlight~1, handing the propulsion stack and the launch authorisation to a single EPC would put flight-critical intellectual property and statutory accountability in a contractor who does not hold the licence \citep{gilmour2025,asa2025}. Two Turner tests decide the model \citep{turner2014}. Scope certainty is high on site operations---seventeen Level-3 packages, locked dates, a \PEPbase\ work estimate---and low on the vehicle. Client capability is the other test: Gilmour can integrate a launch vehicle; it should not pour pad grout or run a 100~t crawler as a core competence. The charter supplied the 40/30/30 weights; A1 did not score the options. This PEP scores them for the site-ops choice only: Strategic Fit 40\% (EPC 2/5, D\&C 5/5), Technical Feasibility 30\% (3/5 vs 4/5), Schedule 30\% (3/5 vs 4/5), weighted 2.6 versus 4.4, so D\&C wins. Alliance contracting was the second rejected alternative: a four-month campaign cannot carry pain-share overhead, and GBRMPA and CASA duties are not commercially shareable \citep{gbrmpa2019,casa2024}.

\subsection{Quality}
\label{sec:delivery-quality}

Quality is specified against AS9100D and ISO~9001:2015 \citep{sae2016,iso9001}. Four hold points, each tied to a register row, are listed in Table~\ref{tab:holdpoints}. HP-1 is pad tie-down weld NDT on A-123 (R-08). HP-2 is the GSE helium leak test on A-124 (R-05). HP-3 is the pre-static-fire ignition interlock on A-133 (R-01). HP-4 is LRR sign-off on A-141: Launch Director, RSO and Juru Officer review with Abdul and Ziyad; the payload sponsor witnesses and does not sign. Work does not proceed past a hold without Ziyad's release or a Board waiver, and a Board waiver is outside Abdul's DoA.

\begin{table}[H]
\centering
\small
\caption[Quality hold points]{Quality hold points for the Bowen campaign. Each hold is a named artefact, not a meeting. Measurement rules in Section~\ref{sec:delivery-evm} follow the same IDs.}
\label{tab:holdpoints}
\begin{tabularx}{\textwidth}{@{}l l l >{\raggedright\arraybackslash}p{3.55cm} >{\raggedright\arraybackslash}X@{}}
\toprule
Hold & Activity & Risk & Artefact & Release \\
\midrule
HP-1 & A-123 (15~Oct) & R-08 & Pad tie-down weld NDT & D\&C mgr / Ziyad \\
HP-2 & A-124 (12~Oct) & R-05 & GSE helium leak test & GSE Lead / Ziyad \\
HP-3 & A-133 (30~Oct) & R-01 & Pre-static-fire ignition interlock & GSE Lead / Launch Director / Ziyad \\
HP-4 & A-141 (10~Nov) & R-16 & LRR sign-off (payload sponsor witnesses) & Launch Director / RSO / Juru Officer / Abdul / Ziyad \\
\bottomrule
\end{tabularx}
\end{table}


Measurement follows the holds and is locked in YAML. Discrete tests and flight artefacts---A-132, A-133, A-141, A-143---are 0/100. A-113 is elapsed-day percent-complete because the work is a 76-day wait, not a binary test. Installations---A-123 and A-124---are milestone percent-complete tied to HP-1 and HP-2. LRR waits on A-134 rather than on A-133: HP-4 cannot be signed on incomplete post-fire data, which is why Option~A was rejected \citep{sae2016}.

\subsection{Health, safety and environment}
\label{sec:delivery-hse}

HSE is this pad under the \textit{Work Health and Safety Act 2011} (Qld) \citep{whs2011}. Liquid oxygen is the process hazard: LOX PPE, oxygen-clean tools, and exclusion during load and detank on A-132 and A-133. This PEP's pad controls---not A1 charter numbers---are a 1.5~km blast zone and twelve simultaneous heads, the same cap on Figure~\ref{fig:resources}. Acoustic treatment is a water deluge so static-fire level at Bowen township stays at or below 135~dB (R-17), with community notice fourteen days prior. Juru artefact buffers and stop-work authority (R-04) are HSE controls: A-114 feeds A-132 as a cultural hold. Finish-line HSE remains as chartered: zero lost-time injuries, zero GBRMPA breaches, and zero Juru heritage breaches, with a five-second static fire of the four first-stage hybrid engines and handover of at least 98.5\% of telemetry. Marine-mammal protocol on A-142 (R-15) is owned by the RSO \citep{gbrmpa2019}. If R-14 fires, Option~C shortens A-134 only and keeps A-141 at five days and A-142 at four days so HP-4 and the marine hold are not crashed. Occupancy never exceeds twelve on the mitigated histogram: R-13 is already treated in the \PEPbase\ base by moving A-125 and hiring four surge technicians.

\subsection{Earned value management}
\label{sec:delivery-evm}

Earned value is the instrument that will tell the Board whether the \PEPbac\ authorisation is still intact \citep{pmi2019evm,fleming2016,kerzner2017}. It is a framework. The campaign has not started, so there is no fourteen-week EV/AC history to plot.

Distributed planned value is \PEPbase, identical to Table~\ref{tab:wbs-cost} and to the S-curve in Figure~\ref{fig:scurve}. Contingency of \PEPcontingency\ is undistributed budget: it is not December cash and it is not monthly PV. Authorised $BAC$ is \PEPbac. Option~C's \PEPoptionC\ is not in the work PMB; if R-14 fires and the Board invokes it, the draw appears as actual cost charged to contingency.

$SV$ and $CV$ are currency, never days:
\begin{equation}
SV = EV - PV, \qquad CV = EV - AC, \qquad SPI = EV/PV, \qquad CPI = EV/AC.
\label{eq:evm}
\end{equation}
In-flight estimate at completion on the work baseline is
\begin{equation}
EAC = AC + (BAC_{\text{work}} - EV)/CPI,
\label{eq:eac}
\end{equation}
with $BAC_{\text{work}} = \PEPbase$. The same CPI is applied to the authorisation ceiling as $BAC \times AC / EV$, reported against \PEPbac \citep{pmi2019evm}. Finish-line rule: final $AC \le \PEPbac$. A Green CPI of \PEPonPlanCPI\ in October is allowed; finishing above $BAC$ is not.

Thresholds that trigger action, used only in this section, are:

\begin{itemize}[leftmargin=1.4em,itemsep=0.25em]
\item \textbf{Green} --- $\PEPgreenLo \le CPI,SPI \le \PEPgreenHi$: dashboard only.
\item \textbf{Amber} --- $\PEPamberLo \le$ index $< \PEPgreenLo$: Exception Report and a seven-day recovery plan within 48~hours.
\item \textbf{Red} --- index $< \PEPredLt$, \emph{or} $T$-0 slip $> \PEPtZeroSlip$~days: contingency freeze and Board intervention. $SV$ is not written in days.
\end{itemize}

The only dashboard numbers permitted in this PEP are the two illustrative readings at M-4 (\PEPmFourLong~2026) in Table~\ref{tab:evm-m4} and Figure~\ref{fig:evm}. Planned value at that date is \PEPpvmFour. The planned label is HP-1 signed (A-123 100\%), HP-2 signed, A-113 elapsed-day percent-complete. On-plan books EV equal to PV (\PEPonPlanEV) and AC \PEPonPlanAC, so $CPI = \PEPonPlanEVnum/\PEPonPlanACnum = \PEPonPlanCPI$ and $SPI = \PEPonPlanEVnum/\PEPpvmFournum = \PEPonPlanSPI$---Green. Even so,
\[
EAC = \PEPonPlanACnum \times \PEPbasenum / \PEPonPlanEVnum = \$\PEPonPlanEacnum,
\]
which is identical to $AC + (BAC_{\text{work}}-EV)/(EV/AC)$. Authorisation
\[
BAC \times AC / EV = \PEPbacnum \times \PEPonPlanACnum / \PEPonPlanEVnum = \$\PEPonPlanIeacnum
\]
sits above the \PEPbac ceiling. Green at M-4 still carries a recovery obligation toward the finish-line rule.

The late-stack illustration holds PV at \PEPpvmFour but drops EV to \PEPlateEV and AC to \PEPlateAC because HP-1 is unsigned. A-123 is milestone 0/100 tied to that hold, so the \$218{,}000 control account contributes \$0 of EV---the whole package is withheld, not a partial-hours claim. Then $CPI = \PEPlateEVnum/\PEPlateACnum = \PEPlateCPI$ and $SPI = \PEPlateEVnum/\PEPpvmFournum = \PEPlateSPI$---\PEPlateBand. Ilsa issues the Exception Report within 48~hours and freezes contingency. If R-14 has fired, the seven-day recovery invokes Option~C and crashes A-134 only---not A-141 and not A-142. Work $EAC$ is $\PEPlateACnum \times \PEPbasenum / \PEPlateEVnum = \$\PEPlateEacnum$. $BAC \times AC / EV = \PEPbacnum \times \PEPlateACnum / \PEPlateEVnum = \$\PEPlateIeacnum$. Those two points are a worked example at one physical gate, not a time series and not a claim that October has already been worked.

\begin{table}[H]
\centering
\footnotesize
\setlength{\tabcolsep}{4.0pt}
\caption[Illustrative earned-value status at M-4]{Illustrative earned-value status at M-4 (15~October~2026). The campaign has not been executed. PV is the locked work-PMB value at stacking complete. Indices use the reported two-decimal CPI/SPI from $EV/AC$ and $EV/PV$.}
\label{tab:evm-m4}
\begin{tabularx}{\textwidth}{@{}>{\raggedright\arraybackslash}p{4.55cm} r r r c c X@{}}
\toprule
Status at M-4 & PV & EV & AC & CPI & SPI & Action \\
\midrule
HP-1 signed (A-123 100%), HP-2 signed, A-113 elapsed-day percent-complete & 865{,}578 & --- & --- & --- & --- & --- \\
Illustrative on-plan & 865{,}578 & 865{,}578 & 901{,}000 & 0.96 & 1.00 & Green --- weekly dashboard \\
Illustrative late stack & 865{,}578 & 800{,}178 & 870{,}000 & 0.92 & 0.92 & Amber --- 48~h exception report \\
\bottomrule
\end{tabularx}
\end{table}


\begin{figure}[!htbp]
\centering
\includegraphics[width=0.98\linewidth]{evm_example.pdf}
\caption[Illustrative earned-value chart at M-4]{Illustrative status at M-4; campaign not yet executed. Upper panel: PV, EV and AC at the single status date \PEPmFourLong~2026 for the on-plan and late-stack cases. Lower panel: Green / Amber / Red CPI--SPI bands, with the on-plan pair (\PEPonPlanCPI, \PEPonPlanSPI) in Green and the late-stack pair (\PEPlateCPI, \PEPlateSPI) in \PEPlateBand. Red is also triggered by $T$-0 slip greater than \PEPtZeroSlip days. No weekly EV/AC history is implied.}
\label{fig:evm}
\end{figure}
